\documentclass[10pt,a4paper]{amsart}
\usepackage[margin=19mm]{geometry}
\usepackage{amsmath,amssymb,mathtools}
\usepackage[scr=rsfs,cal=euler]{mathalfa}
\usepackage{xfrac}
\usepackage[unicode,hidelinks,bookmarks=false]{hyperref}
\newcommand{\ie}{i.e.\ }
\newcommand{\cf}{cf.\ }
\newcommand{\resp}{respectively\ }
\newcommand{\Subsection}{\subsection}
\providecommand{\nobreakdash}{\nobreak\hbox{-}}
\setlength{\emergencystretch}{2em}
\multlinegap0pt
\usepackage{bm}
\usepackage{bbm}
\usepackage{dsfont}
\usepackage{xspace}
\usepackage{cancel}
\usepackage{mathtools}
\usepackage{amsthm}
\makeatletter
\@ifundefined{lemma}{\newtheorem{lemma}{Lemma}}{}
\makeatother
\makeatletter
\expandafter\ifx\csname proof\endcsname\relax
\newenvironment{proof}[1][Proof]{\textbf{#1.} }{\ \rule{0.5em}{0.5em}}
\fi
\makeatother
\title[On the Spectrality of the Differential Operators with Period]{On the Spectrality of the Differential Operators with Periodic Coefficients}
\author{O. A. Veliev}
\date{Source: arXiv:2504.07873v4 (CC BY 4.0)}
\begin{document}
\maketitle

\noindent\emph{Source and licence.} On the Spectrality of the Differential Operators with Periodic Coefficients. Version arXiv:2504.07873v4 (CC BY 4.0), distributed under the Creative Commons Attribution 4.0 International licence, as recorded in the arXiv licence metadata for this exact version and shown on the abstract page https://arxiv.org/abs/2504.07873v4. Full rights notice: licenses/arxiv-2504.07873-CC-BY-4.0.txt.\par\smallskip

\noindent\textit{Editorial scope.} Counted unit: the Lemma whose source label is \texttt{None} in the article (source0.tex lines 616--625, source label \texttt{None}), delivered together with the article's own complete original proof of it (source0.tex lines 627--816, one \texttt{proof} environment of 190 lines). No mathematics is written, repaired or re-derived: statement and proof are the article's own text line for line. Required context retained uncounted: none. Every definition, display and label of those lines is retained unchanged. Omitted from this package and not redistributed: 1--615, 626--626, 817--1038. Nothing inside a retained excerpt is omitted or altered: every retained body is delivered exactly as the article's lines are written, so no omission receipt of any kind accompanies this package. Citations: the retained excerpts carry no citation command at all, so no bibliography entry of the article is transcribed and no reference is redistributed. Reference adaptation: none was needed and none was made; every reference inside the delivered unit resolves inside it, so no unresolved reference or citation is produced. Printed slips kept literal: no printed word, symbol, hypothesis or number of the counted statement or of its proof was altered. Layout and class: the article's own class is \texttt{article} with packages including graphicx, amsmath, amsfonts, amssymb; the wrapper is the lane's pinned \texttt{amsart} editorial wrapper at 10pt on A4, which restates the theorem-like environments the retained text needs and adds the article's own macro definitions that the retained text needs (none), with extra wrapper packages (bm, bbm, dsfont, xspace, cancel, mathtools). The delivered numbering is the unit's own. Assets: no retained span contains a figure environment, a graphics-inclusion command, a PDF-page inclusion, a table or a TikZ picture; no image file is copied into this package, the package declares no asset dependency and no image file is redistributed with it. Licence: version arXiv:2504.07873v4 of this article is distributed under the Creative Commons Attribution 4.0 International licence, as recorded in the arXiv licence metadata for this exact version and shown on the abstract page https://arxiv.org/abs/2504.07873v4. A full rights notice is delivered at licenses/arxiv-2504.07873-CC-BY-4.0.txt. Characters: every retained excerpt is pure ASCII, so the delivered engine, XeLaTeX with its default Latin Modern text font, reports no missing character.
\begin{lemma}
If $n$ is an even number greater that $2$ and condition (4) is satisfied, then
the horizontal lines
\[
H(n,t,s)=\left\{  (x,y)\in\mathbb{R}^{2}:y=ci^{n-2}((2s+1)\pi+\pi
t)^{n-1}\right\}
\]
for $s=0,\pm1,\pm2,...,$ belong to the resolvent set of $L_{t,\varepsilon}$
for all $\varepsilon\in\lbrack0,1]$ and $t\in\lbrack0,1].$
\end{lemma}

\begin{proof}
\bigskip First, consider the case $s\geq0$. Suppose there exists $\lambda\in
H(n,t,s)$ which is an eigenvalue of $L_{t,\varepsilon}$ for some
$\varepsilon\in\lbrack0,1].$ Then, for the denominator in the expression from
(17), we have the estimate:
\begin{equation}
\left\vert \left(  \lambda-\left(  2\pi k+\pi t\right)  ^{n}-c(2\pi ki+\pi
ti)^{n-1}\right)  \right\vert \geq\left\vert c((2s+1)\pi+\pi t)^{n-1}-c(2\pi
k+\pi t)^{n-1}\right\vert .\tag{18}%
\end{equation}
Using this estimation we prove that
\begin{equation}
\sum_{k\in\mathbb{Z}}\left\vert \left(  \Psi_{\lambda},e^{i\left(  2\pi k+\pi
t\right)  x}\right)  \right\vert ^{2}<1.\tag{19}%
\end{equation}
This contradicts Parseval's equality for the orthonormal basis $\left\{
e^{i\left(  2\pi k+\pi t\right)  x}:k\in\mathbb{Z}\right\}  .$ This means that
$\lambda$ is not an eigenvalue and therefore belong to the resolvent set of
the operators $L_{t,\varepsilon}$. To prove (19), we write the left-hand side
of (19) as the sum of the following four terms: $S_{1}(s),$ $S_{2}(s),$
$S_{2}(s),$ $S_{4}(s),$ and estimate them separately, where
\[
S_{1}(s)=\sum_{k>s}\left\vert \left(  \Psi_{\lambda},e^{i\left(  2\pi k+\pi
t\right)  x}\right)  \right\vert ^{2},\text{ }S_{2}(s)=\sum_{-s\leq k\leq
s,k\neq0}\left\vert \left(  \Psi_{\lambda},e^{i\left(  2\pi k+\pi t\right)
x}\right)  \right\vert ^{2},
\]%
\[
S_{3}(s)=\sum_{k<-s}\left\vert \left(  \Psi_{\lambda},e^{i\left(  2\pi k+\pi
t\right)  x}\right)  \right\vert ^{2},\text{ }S_{4}(s)\left\vert =\left(
\Psi_{\lambda},e^{i\pi tx}\right)  \right\vert ^{2}.
\]
To estimate $S_{1}(s)$ and $S_{2}(s),$ we use the following obvious
inequalities
\begin{equation}
\left\vert a^{n-1}-b^{n-1}\right\vert \geq\left\vert a-b\right\vert \left\vert
a\right\vert ^{n-2},\tag{20}%
\end{equation}
for $ab\geq0$ and for $\left\vert a\right\vert \leq b,$ respectively. If
$k>s\geq0,$ then both $a=:(2\pi k+\pi t)$ and $b:=((2s+1)\pi+\pi t)$ are
positive numbers, if $s\geq0$ and $-s\leq k\leq s,$ then $\left\vert
a\right\vert \leq b.$ Therefore, using (18) and (20), we obtain
\[
\left\vert \left(  \lambda-\left(  2\pi k+\pi t\right)  ^{n}-c(2\pi ki+\pi
ti)^{n-1}\right)  \right\vert \geq\left\vert c\pi(2k-2s-1))(2\pi k+\pi
t)^{n-2}\right\vert .
\]
Substituting this into (17), where $P(k,t)=\left(  2\pi k+\pi t\right)
^{n-2}$ for $k\neq0$ and using the well-known identity
\begin{equation}
\sum_{n=1}^{\infty}\frac{1}{(2n-1)^{2}}=\frac{1}{8}\pi^{2},\tag{21}%
\end{equation}
we obtain the following estimates
\begin{equation}
S_{1}(s)\leq\sum_{k>s}\frac{C^{2}}{c^{2}\pi^{2}\left\vert (2k-2s-1)\right\vert
^{2}}=\frac{C^{2}}{8c^{2}},\text{ }S_{2}(s)<\frac{C^{2}}{8c^{2}}.\tag{22}%
\end{equation}


If $k<0,$ then both $((2s+1)\pi+\pi t)^{n-1}$ and $-(2\pi k+\pi t)^{n-1}$ are
positive numbers for $s\geq0.$ Therefore, from (18), we obtain
\begin{equation}
\left\vert \left(  \lambda-\left(  2\pi k+\pi t\right)  ^{n}-c(2\pi ki+\pi
ti)^{n-1}\right)  \right\vert \geq\left\vert c\right\vert \left\vert (2\pi
k+\pi t)^{n-1}\right\vert . \tag{23}%
\end{equation}
Moreover, $\left\vert 2\pi k+\pi t\right\vert \geq\pi\left(  \left\vert
2k\right\vert -1\right)  $ for $t\in\lbrack0,1].$ Using this inequality (17),
(23) and (21), we obtain
\begin{equation}
S_{3}(s)\leq\sum_{k<-s}\frac{C^{2}}{c^{2}\pi^{2}(\left\vert 2k\right\vert
-1)^{2}}=\frac{C^{2}}{c^{2}\pi^{2}}\left(  \frac{\pi^{2}}{8}-\sum_{-s\leq
k<0}\frac{1}{(\left\vert 2k\right\vert -1)^{2}}\right)  . \tag{24}%
\end{equation}
It remains to estimate $S_{4}(s).$ Using (17), where $P(0,t)=\pi^{n-2},$ and
(18), we obtain
\begin{equation}
S_{4}(s)\leq\frac{\left(  \pi^{n-2}C\right)  ^{2}}{c^{2}\left\vert
((2s+1)\pi+\pi t)^{n-1}-(\pi t)^{n-1}\right\vert ^{2}}. \tag{25}%
\end{equation}
Now, using (22), (24) and (25), we prove that%
\begin{equation}
\sum_{k\in\mathbb{Z}}\left\vert \left(  \Psi_{\lambda},e^{i\left(  2\pi k+\pi
t\right)  x}\right)  \right\vert ^{2}=S_{1}(s)+S_{2}(s)+S_{3}(s)+S_{4}%
(s)<(\frac{1}{4}+\frac{1}{\pi^{2}})\frac{C^{2}}{c^{2}}. \tag{26}%
\end{equation}
for all $s\geq0.$ First we prove (26) for $s=0.$ From (24) and (25) , we
obtain that
\begin{equation}
S_{3}(0)\leq\frac{C^{2}}{8c^{2}},\text{ }S_{4}(0)\leq\frac{\left(  \pi
^{n-2}C\right)  ^{2}}{c^{2}\left\vert (\pi+\pi t)^{n-1}-(\pi t)^{n-1}%
\right\vert ^{2}}\leq\frac{C^{2}}{\pi^{2}c^{2}}, \tag{27}%
\end{equation}
since $(\pi+\pi t)^{n-1}-(\pi t)^{n-1}$ is a nondecreasing function on
$[0,1].$ Moreover, it follows from the definition of $S_{2}(s)$ that
$S_{2}(0)=0.$ Therefore, the inequality in (26), for $s=0,$ follows from (22)
and (27). Now, we prove (26) for $s\geq1.$ It follows from (24) and (25) that
\begin{equation}
S_{3}(s)\leq\frac{C^{2}}{c^{2}\pi^{2}}\left(  \frac{\pi^{2}}{8}-1\right)
,\text{ }S_{4}(s)\leq\frac{\left(  \pi^{n-2}C\right)  ^{2}}{c^{2}\left\vert
(3\pi+\pi t)^{n-1}-(\pi t)^{n-1}\right\vert ^{2}}\leq\frac{1}{9}\frac{C^{2}%
}{\pi^{2}c^{2}} \tag{28}%
\end{equation}
for all $s\geq1$ and $n\geq2,$ since $(3\pi+\pi t)^{n-1}-(\pi t)^{n-1}$ is an
increasing function on $[0,1].$ Instead (27) using (28) and noting that
\[
\frac{1}{\pi^{2}}+\frac{8}{9}\frac{1}{\pi^{2}}>\frac{1}{8},
\]
we see that (26) holds for all $s\geq1.$

Now let us consider the case $s<0.$ If $k<0,$ then both $(2\pi k+\pi t)$ and
$(2s+1)\pi+\pi t$ are nonpositive numbers and we can use (20). Therefore,
arguing as in the proof of (22), we obtain
\begin{equation}
S_{5}(s):=\sum_{k\leq s}\left\vert \left(  \Psi_{\lambda},e^{i\left(  2\pi
k+\pi t\right)  x}\right)  \right\vert ^{2}\leq\frac{C^{2}}{8c^{2}}\text{ ,
}S_{6}(s):=\sum_{s<k<0}\left\vert \left(  \Psi_{\lambda},e^{i\left(  2\pi
k+\pi t\right)  x}\right)  \right\vert ^{2}<\frac{C^{2}}{8c^{2}}.\tag{29}%
\end{equation}
If $k>0,$ then both $((2s+1)\pi+\pi t)^{n-1}$ and $-(2\pi k+\pi t)^{n-1}$ are
nonpositive numbers for $s<0.$ Therefore, from (18) and (17), by using the
inequality $\left\vert 2\pi k+\pi t\right\vert \geq\left\vert 2\pi
k\right\vert $ for $t\in\lbrack0,1]$ and the well-known equality
\[
\sum_{n=1}^{\infty}\frac{1}{(2n)^{2}}=\frac{1}{24}\pi^{2},
\]
we obtain%
\[
\left\vert \left(  \lambda-\left(  2\pi k+\pi t\right)  ^{n}-c(2\pi ki+\pi
ti)^{n-1}\right)  \right\vert \geq\left\vert c(2\pi k+\pi t)^{n-1}\right\vert
\]
and
\begin{equation}
S_{7}(s):=\sum_{k>0}\left\vert \left(  \Psi_{\lambda},e^{i\left(  2\pi k+\pi
t\right)  x}\right)  \right\vert ^{2}\leq\sum_{k>0}\frac{C^{2}}{c^{2}\pi
^{2}(2k)^{2}}=\frac{C^{2}}{24c^{2}}.\tag{30}%
\end{equation}


It remains to estimate $\left\vert \left(  \Psi_{\lambda},e^{i\pi tx}\right)
\right\vert $ for $s<0.$ First consider the case $s=-1.$ Since the expression
$\left\vert (-\pi+\pi t)^{n-1}-(\pi t)^{n-1}\right\vert $ gets its minimum
value at $t=\frac{1}{2},$ we have
\[
\left\vert (-\pi+\pi t)^{n-1}-(\pi t)^{n-1}\right\vert \geq2^{2-n}\pi^{n-1}%
\]
for $s=-1$ and $t\in\lbrack0,1].$ Now, from (18) and (17), where
$P(0,t)=\pi^{n-2},$ we obtain, \
\begin{equation}
\left\vert \left(  \Psi_{\lambda},e^{i\pi tx}\right)  \right\vert ^{2}%
\leq\frac{2^{2n-4}C^{2}}{\pi^{2}c^{2}}. \tag{31}%
\end{equation}
Therefore using (29)-(31) and noting that $S_{6}(-1)=0,$ we see that
\begin{equation}
\sum_{k\in\mathbb{Z}}\left\vert \left(  \Psi_{\lambda},e^{i\left(  2\pi k+\pi
t\right)  x}\right)  \right\vert ^{2}<\left(  \frac{1}{6}+\frac{2^{2n-4}}%
{\pi^{2}}\right)  \frac{C^{2}}{c^{2}} \tag{32}%
\end{equation}
for $s=-1.$

If $s<-1,$ then
\[
\left\vert ((2s+1)\pi+\pi t)^{n-1}-(\pi t)^{n-1}\right\vert \geq\left\vert
(-3\pi+\pi t)^{n-1}-(\pi t)^{n-1}\right\vert \geq(2^{n-1}+1)\pi^{n-1}%
\]
and \
\begin{equation}
\left\vert \left(  \Psi_{\lambda},e^{i\pi tx}\right)  \right\vert ^{2}%
\leq\left(  \frac{1}{(2^{n-1}+1)^{2}\pi^{2}}\right)  \frac{C^{2}}{c^{2}}.
\tag{33}%
\end{equation}
This with (29) and (30) implies that
\begin{equation}
\sum_{k\in\mathbb{Z}}\left\vert \left(  \Psi_{\lambda},e^{i\left(  2\pi k+\pi
t\right)  x}\right)  \right\vert ^{2}<\left(  \frac{7}{24}+\frac{1}%
{(2^{n-1}+1)^{2}\pi^{2}}\right)  \frac{C^{2}}{c^{2}} \tag{34}%
\end{equation}
for $s<-1.$ Thus, by (26), (32) and (34) we have
\[
\sum_{k\in\mathbb{Z}}\left\vert \left(  \Psi_{\lambda},e^{i\left(  2\pi k+\pi
t\right)  x}\right)  \right\vert ^{2}<A\frac{C^{2}}{c^{2}},
\]
for all $\lambda\in H(n,t,s),$ $s\in\mathbb{Z}$ and $t\in\lbrack0,1],$ where
\[
A=\max\left\{  \frac{1}{4}+\frac{1}{\pi^{2}},\frac{1}{6}+\frac{2^{2n-4}}%
{\pi^{2}},\frac{7}{24}+\frac{1}{(2^{n-1}+1)^{2}\pi^{2}}\right\}  =\frac{1}%
{6}+\frac{2^{2n-4}}{\pi^{2}}%
\]
for $n\geq4.$ It means that if (4) holds then (19) is satisfied. The lemma is proved.
\end{proof}

\end{document}
